\section{Conclusion}
\label{sec:conclusion}

PatchTST gives every input patch the same trust and normalizes away the level and scale of each window, which makes it sensitive to spikes and to shifts in the input. We proposed RA-PatchTST, which adds a patch reliability gate, two shift features for the forecasting head and corruption-aware training, at a cost of 993 parameters, and evaluated it on ETTh1 against a PatchTST re-implementation that reproduces the published clean error to within 1.1\%. Across 21 corrupted test conditions the mean MSE fell from 0.429 to 0.413 ($-3.7\%$), with reductions of 11.0\% under shift, 6.5\% under spikes and 5.2\% under mixed damage, and of 20.1\% at the most severe shift; the clean error rose by 0.7\%.

The ablation qualifies these numbers. Training on corrupted inputs provides about four fifths of the gain. The shift-aware head adds most of the remainder and the gate a little more, and the two effects are additive. The gate discounts spiky patches but does not identify missing values, and its weights stay close to one. The study covers one dataset, one horizon and two seeds, with synthetic corruptions from the training families, so the evidence is preliminary. Within those limits, the results suggest that for patch-based forecasters a large part of the robustness to corrupted inputs can be bought cheaply, through the training data and a global view of shift, and that learning to distrust individual patches needs a stronger signal than the forecasting loss provides.
